% !TeX root = ../main.tex
\chapter{方法与排版示例}[Methods and Typesetting Examples]
\section{数学公式}[Equations]
公式~\ref{eq:mean} 用于检验求和字体与按章编号。
\begin{equation}
  \bar x = \frac{1}{N}\sum_{i=1}^{N}x_i .
  \label{eq:mean}
\end{equation}
\section{算法}[Algorithm]
算法~\ref{alg:mean} 保留中期模板的中文输入输出、五号字、行号与块竖线。
\begin{algorithm}[htbp]
\caption{样本均值计算示例}\label{alg:mean}
\KwIn{样本 $x_1,\ldots,x_N$，样本数 $N\ge1$}
\KwOut{样本均值 $\bar x$}
$s\gets0$\;
\For{$i=1,\ldots,N$}{
  $s\gets s+x_i$\;
}
$\bar x\gets s/N$\;
\end{algorithm}
\section{图与表}[Figures and Tables]
图~\ref{fig:placeholder} 和表~\ref{tab:example} 仅为排版示例，提交前应替换。
\begin{figure}[htbp]
\centering
\fbox{\parbox[c][20mm][c]{0.65\textwidth}{\centering 替换为研究图片}}
\bicaption[fig:placeholder]{}{图片占位示例}{Fig.}{Figure placeholder}
\end{figure}
\begin{table}[htbp]
\centering
\bicaption[tab:example]{}{三线表示例}{Table}{Example table}
\begin{tabular}{lc}
\toprule
项目 & 示例值 \\
\midrule
样本数 & 100 \\
\bottomrule
\end{tabular}
\end{table}
